\documentclass[11pt,a4paper]{article}
\usepackage[margin=19mm,headheight=15pt]{geometry}
\usepackage{fontspec}
\setmainfont{Noto Serif}
\usepackage{amsmath,amssymb,mathtools}
\usepackage{xcolor}
\usepackage{fancyhdr}
\usepackage{enumitem}
\usepackage{microtype}
\setlength{\parindent}{0pt}
\setlength{\parskip}{6pt}
\definecolor{navy}{HTML}{183153}
\definecolor{lightgray}{HTML}{F3F5F7}
\pagestyle{fancy}\fancyhf{}
\lhead{\small ME 5309 -- Mid-Semester Examination}
\rhead{\small ME 5309}
\cfoot{\thepage}
\begin{document}
\begin{center}
{\Large\bfseries ME 5309 -- Mid-Semester Examination}\\[3pt]
{\large Answer Booklet}\\[8pt]
\end{center}
\textbf{Instructions:} Attempt any three questions. Each question carries 10 marks. Show working where appropriate.
\section*{Question 1 \hfill [10 marks]}
\subsection*{Q1(a) \hfill [4 marks]}

(i) f\_x = 16/5 = 3.2 pixels, using focal length 16 mm and pixel size 5 microm.
(ii) Sensor width is 6.4 mm, so field width = 400x6.4/16 = 160 mm and scale = 0.125 mm/pixel.

\subsection*{Q1(b) \hfill [3 marks]}

\vspace{2mm}

\subsection*{Q1(c) \hfill [3 marks]}

For zero-order hold, F = e\textasciicircum{}(AT) and G = integral\_0\textasciicircum{}T e\textasciicircum{}(Aτ)B dτ.
Here:\[\begin{aligned}A = -2, B = 1, T = 0.1 s\end{aligned}\]
Thus:\[\begin{aligned}F = e^(-0.2) = 0.818731\\ G = (1-e^(-0.2))/2 = 0.0906346\end{aligned}\]
Therefore:\[\begin{aligned}x_{k+1} = 0.818731 x_k + 0.0906346 u_k\end{aligned}\]\par\smallskip With $AT=(-2)(0.1)=-0.2$, $F=e^{-0.2}=0.8187308$. For $G$, $(1-F)/2=(1-0.8187308)/2=0.1812692/2=0.0906346$. Thus the input contribution is positive.

\section*{Question 2 \hfill [10 marks]}
\subsection*{Q2(a) \hfill [2 marks]}

ẋ₁ = x₂ and ẋ₂ = (τ-b x₂+mgℓ sin θ)/J.
At equilibrium:\[\begin{aligned}τ_0 = -mg\ell sin \theta_e,\\ \text{an equilibrium exists for |τ₀| less than or equal to 2.4525 N m}\end{aligned}\]

\subsection*{Q2(b) \hfill [2 marks]}

\[\begin{aligned}A = [[0,1], [-mg\ell sin(60^{\circ})/J, -b/J]] = [[0,1], [-106.2, -2.5]],\\ B = [[0],[50]]\end{aligned}\]

\subsection*{Q2(c) \hfill [4 marks]}

At steady state θ̇ = 0, and I assume θss = θd = 90° because the PD controller tracks the desired angle.
Thus the steady-state error is zero.

\subsection*{Q2(d) \hfill [2 marks]}

The gravity term changes the equilibrium, so the steady-state error is still nonzero.
The poles should change.

\section*{Question 4 \hfill [10 marks]}
\subsection*{Q4(a) \hfill [2 marks]}

Closed-loop characteristic polynomial is s²+(a+bk₂)s+(c+bk₁)=s²+4s+8, giving poles -2±2j.
At DC, θ/r = bNff/(c+bk₁), so unity gain requires Nff=(c+bk₁)/b=(6+2)/1=8.\par\smallskip Comparing coefficients with $s^2+4s+8$ gives $a+bk_2=4$ and $c+bk_1=8$. The feedforward gain is chosen from the nominal plant so the reference scaling is unity; here this gives $N_{ff}=8$. The closed-loop roots of $s^2+4s+8$ are $-2\pm2j$.

\subsection*{Q4(b) \hfill [2 marks]}

\[\begin{aligned}\theta_{ss} = 0.75\times8/(6+2)=0.75,\\ \text{the percentage tracking error is 25\%}\end{aligned}\]

\subsection*{Q4(c) \hfill [4 marks]}

The desired polynomial is (s+4)(s²+4s+8).
The closed-loop polynomial should be cubic because an integrator is added.

\subsection*{Q4(d) \hfill [2 marks]}

\vspace{2mm}

\end{document}